\chapter{Session Preparation Guide} \index{Preparation Guide}

\newthought{Each session is prepared by one student group.} This chapter
explains what the preparing group is responsible for, how to submit
improvements, and what is expected on the day of the session.

\vspace{1em}


\noindent\begin{tabular}{@{}p{3.8cm}p{6.5cm}@{}}
    \textbf{Phase} & \textbf{What you do} \\[4pt]
    Before the session &
        Review the lecture notes, propose improvements via pull request,
        and prepare the debate (including providing the moderator). \\[6pt]
    During the session &
        Moderate the debate, introduce any additions you made, and
        facilitate the wrap-up discussion. \\
\end{tabular}

% ---------------------------------------------------------------
\newthought{Read the chapter critically.} Go through every section of the
assigned session chapter: the guided reading sheet, the theoretical margin
notes, the debate roles and positions, and the key learnings. You own that Chapter now, you can touch everything. Ask
yourselves:

\begin{itemize}
    \item Are the guided reading questions clear, fair, and thought-provoking?
          Could any of them be sharpened, split, or replaced?
    \item Is there a question missing that would strengthen the discussion?
    \item Are the theoretical concepts in the margin notes accurate and
          up to date? Are there newer or better sources, or extending ones?
    \item Could a figure, diagram, image, or short quote make a concept
          more accessible or more memorable?
    \item Are the debate positions balanced? Do the talking points give
          each role enough material to argue convincingly? Is there a notion not captured in the roles?
\end{itemize}

\newthought{Submit your improvements as a pull request.} All lecture notes
are maintained in a shared repository. The preparing group submits their
proposed changes as a \textbf{pull request (PR)} directly on the
repository. This ensures that improvements are versioned, reviewable,
and attributable. Make sure to split PRs by the categories provided in the next section,
and to provide a clear description of each change and its rationale. The PR must compile without errors. 
The instructor reviews and merges the PR after the session.

\subsection*{What belongs in a pull request}

\begin{enumerate}
    \item \textbf{Question edits.} Rephrase, split, or replace a guided
          reading question. Add a new question if you identify a gap.
          Briefly explain in the PR description why the change improves the
          sheet.

    \item \textbf{Source updates.} Add a recent paper, report, or case
          study that strengthens a theoretical concept or a debate position.
          Add the entry to \texttt{philosophyofai\_references.bib} and cite
          it with \texttt{\textbackslash cite\{\}} in the appropriate margin
          note or body text.

    \item \textbf{Margin notes.} Add a short margin note
          (\texttt{\textbackslash marginnote\{\}}) that provides context,
          a definition, a counter-example, or a link to current events.
          Keep margin notes concise: two to four sentences.

    \item \textbf{Figures and images.} Add a figure to the
          \texttt{figures/} directory and include it with
          \texttt{\textbackslash includegraphics}. Every figure must have a
          caption that explains what it shows and why it matters. Cite the
          source in the caption or in a margin note.

    \item \textbf{Debate refinements.} Strengthen a talking point, add a
          new one, or rebalance the positions if you believe one side is
          under-argued.
\end{enumerate}

    % ---------------------------------------------------------------
\section*{Before the Session: Preparing the Debate}

\newthought{The preparing group provides the moderator.} One member of the
group serves as moderator for the session. The moderator's duties are
described in detail in the Debate Format \& Moderator Guide (Appendix~A).

\newthought{The rest of the group} is ready to step in if the moderator needs assistance, or replacement.
